\section{Complete certification of the fixed cohort}
\label{sec:proof-closure}
A targeted follow-up closes all 16 intervals left open by the earlier
postprocessing. Every one of the 117 instances now has a valid labeling
attaining an independently verified lower bound. Ten retained upper bounds
improve. The disjoint evidence categories are 91 theory/clique certificates,
15 weighted-capacity certificates, and 11 clausal UNSAT proofs. These are
post-hoc, instance-specific results selected after inspecting difficult
cases. They do not change the original timed SAT--ILP observations or
establish a general performance advantage.

\subsection{A weighted lower bound}
Assign nonnegative integer weights $w_v$, not all zero, to the vertices of
$G^2$, and let
\[
 W=\sum_v w_v,\qquad C=\max_{I\text{ independent in }G^2}\sum_{v\in I}w_v.
\]
For $h\ge k\ge1$,
\begin{equation}
 \lambda_{h,k}(G)\ge k\bigl(\lceil W/C\rceil-1\bigr).
 \label{eq:weighted-capacity}
\end{equation}
Indeed, vertices whose labels belong to a block of $k$ consecutive integers
form an independent set of $G^2$. A labeling of span $s$, normalized to
minimum zero, therefore satisfies
$W\le C(\lfloor s/k\rfloor+1)$, proving
\eqref{eq:weighted-capacity}. This is the classical fractional-coloring
dual bound, rounded to an integer number of colors, rather than a new
general theorem. Numerical optimization proposes weights only: the checker
uses integer arithmetic to enumerate all maximal independent sets and
recompute $C$. Nonnegative weights justify restricting the enumeration to
maximal sets. The six difficult graph squares have 170--390 such sets;
this favorable size does not imply polynomial complexity in general.

\subsection{Witness constructions}
The earlier class-order construction extends to arbitrary independent
classes of $G^2$. For $k\le h\le2k$, successive classes receive gap $h$
when they contain an original edge and gap $k$ otherwise. Nonconsecutive
classes differ by at least $2k\ge h$, proving validity. Searching for an
order is an upper-bound procedure, not a complete labeling decision method.
Parameter transfer supplies further witnesses: multiplying an $L(a,b)$
labeling by an integer $t>0$ produces an $L(h,k)$ labeling whenever
$ta\ge h$ and $tb\ge k$. In particular, doubling an $L(2,1)$ witness
gives an $L(3,2)$ witness. These constructions, an additional SAT witness,
and the capacity bounds close 14 of the 16 previously open intervals.

\subsection{Color, rank and residual formulation}
The last two intervals require infeasibility proofs. At proposed span $s$,
put $q=\lfloor s/k\rfloor+1$ and represent labels by
\[
 f(v)=k\pi(c(v))+r_v,\qquad
 c(v)\in\{0,\ldots,q-1\},\quad \pi\in S_q,\quad 0\le r_v<k.
\]
The internal colors $c$ form a proper coloring of $G^2$; $\pi$ assigns
their physical ranks. Every feasible labeling admits this representation
by division with remainder. A square clique may be fixed to distinct
internal color names because $\pi$ remains free. This does not fix its
vertices' numerical labels. One-hot constraints use the sequential
cardinality encoding of Sinz~\cite{sinz2005cardinality}.

For $L(2,1)$, residuals vanish and classes containing an original edge
cannot occupy consecutive ranks. For $L(3,2)$, each vertex retains a
residual bit. Across consecutive ranks, an original edge requires residuals
$(0,1)$ in increasing-rank order; a distance-two pair forbids $(1,0)$.
Rank differences of at least two imply label difference at least three.
The final rank's residual is restricted by $f(v)\le s$.
These conditions are necessary and sufficient, so the formulas are
equisatisfiable with the bounded-span labeling problem. Omitting residual
bits would invalidate this equivalence. Exhaustive tests compare both
specializations with independent label enumeration on every graph with
at most four vertices, across spans 0--7.

For ER\_40\_p0.15\_seed0, the formulas are UNSAT at span 13 for
$L(2,1)$ and at span 27 for $L(3,2)$. Valid labelings of spans 14 and 28
therefore establish the exact optima. Nine additional instances previously
dependent on solver lower bounds receive proofs from the original order
formulation. Two other such instances are certified by the capacity bound.

\subsection{Independent proof replay}
Glucose produces clausal refutations. DRAT-trim checks them and exports
RUP-only LRAT derivations with propagation hints~\cite{cruzfilipe2017lrat}.
A separate standard-library checker reconstructs each CNF from the archived
graph and validates every indicated unit-propagation chain, requiring a
derived empty clause. All 11 proofs pass. The largest uses 57,877 derived
clauses and 5,075,028 checked hints. The checker is tested but has not itself
been formally verified in a proof assistant.

\begin{table}[H]
\centering\small
\caption{Exact optima for the six graphs with formerly open intervals.}
\begin{tabular}{lrrr}
\toprule
Graph & $\lambda_{1,1}$ & $\lambda_{2,1}$ & $\lambda_{3,2}$ \\
\midrule
ER\_40\_p0.15\_seed0 & 13 & 14 & 28 \\
ER\_40\_p0.15\_seed1 & 13 & 13 & 26 \\
ER\_40\_p0.15\_seed2 & 13 & 13 & 26 \\
ER\_60\_p0.15\_seed0 & 22 & 22 & 44 \\
ER\_60\_p0.15\_seed1 & 22 & 22 & 44 \\
ER\_60\_p0.15\_seed2 & 20 & 20 & 40 \\
\bottomrule
\end{tabular}

\tablenote{All graphs are ER with $p=0.15$. Each value has a validated
labeling and a separately checked lower bound. These are supplementary
certificates, not additional successes within the original time limit.}
\end{table}
